\chapter{Higher-order accuracy stencils}
\label{chap:dev:higher-order-accuracy-stencils}

\begin{quote}
Phase metadata. Spec dir:
\passthrough{\lstinline!specs/2026-05-03-phase-7-higher-order-stencils/!}.
Adds \passthrough{\lstinline!stencil::Coefficients<4>!},
\passthrough{\lstinline!stencil::FirstDerivative<4>!}, and templates
\passthrough{\lstinline!diff::laplacian!},
\passthrough{\lstinline!diff::gradient!},
\passthrough{\lstinline!diff::divergence!} on an
\passthrough{\lstinline!Order!} parameter (default
\passthrough{\lstinline!2!}). Version target:
\passthrough{\lstinline!0.7.0!}.
\end{quote}

\subsection{1. Theory}\label{theory}

Central-difference stencils approximate a derivative at \(x_i\) as a
weighted sum of samples at \(x_{i-r}, x_{i-r+1}, \ldots, x_{i+r}\). The
weights \(w_m\) are chosen so the weighted sum matches the target
derivative through some order in the Taylor expansion of \(f\) around
\(x_i\).

Two \textbf{symmetry constraints} organize the family:

\begin{itemize}
\tightlist
\item
  For an \textbf{odd} derivative (\(\partial^k\) with \(k\) odd), the
  stencil is \textbf{anti-symmetric}: \(w_{-m} = -w_m\), so the center
  weight \(w_0\) is exactly zero. This is why
  \passthrough{\lstinline!FirstDerivative<2>::weights[1] = 0!} and
  \passthrough{\lstinline!FirstDerivative<4>::weights[2] = 0!}.
\item
  For an \textbf{even} derivative, the stencil is \textbf{symmetric}:
  \(w_{-m} = +w_m\), with no constraint on \(w_0\).
\end{itemize}

Doubling the half-width \(r\) buys two more matched Taylor orders by
symmetry --- odd-order error terms cancel automatically. So central
differences come in \textbf{even} truncation orders only: 2, 4, 6,
\ldots{} There is no order-3 symmetric central-difference stencil for
these derivatives; an order-3 method would have to be asymmetric
(upwind), which is a different design space.

For the Phase 1 \passthrough{\lstinline!Grid!} (Cartesian-orthogonal
with per-axis cell sizes \(h_x, h_y, h_z\)), the 3D Laplacian decomposes
into independent per-axis 1D second-derivative stencils:

\[
\nabla^2_h \psi_{ijk} \;=\; \sum_{a \in \{x,y,z\}} \frac{1}{h_a^2}\,\sum_{m=-r}^{+r} w_m\,\psi(\mathbf{x}_{ijk} + m\,h_a\,\hat{\mathbf{e}}_a),
\]

with the same weight table \(\{w_m\}\) on every axis (because the
weights are unitless --- they only depend on the integer offsets
\(\{-r, \ldots, +r\}\), not on \(h_a\)). This is what lets one
\passthrough{\lstinline!Coefficients<Order>!} table drive the operator
on \textbf{anisotropic Cartesian-orthogonal} grids without any per-axis
customization.

\subsection{2. Math derivation}\label{math-derivation}

\subsubsection{Order-4 second-derivative
stencil}\label{order-4-second-derivative-stencil}

We want \(\partial^2 f / \partial x^2\) at \(x_i\) from the symmetric
stencil

\[
\partial^2_x f(x_i) \;\approx\; \frac{1}{h^2}\,\sum_{m=-2}^{2} w_m\,f(x_{i+m})
\;=\; \frac{1}{h^2}\bigl[w_{-2} f_{i-2} + w_{-1} f_{i-1} + w_0 f_i + w_1 f_{i+1} + w_2 f_{i+2}\bigr].
\]

Symmetry: \(w_{-m} = w_m\). Three free parameters: \(w_0\),
\(w_1 = w_{-1}\), \(w_2 = w_{-2}\).

Taylor-expand around \(x_i\):

\[
f(x_i + m\,h) \;=\; \sum_{p \ge 0} \frac{(m\,h)^p}{p!}\,f^{(p)}(x_i).
\]

The weighted sum is

\[
\sum_m w_m\,f(x_i + m\,h) \;=\; \sum_{p \ge 0} \frac{h^p}{p!}\,f^{(p)}(x_i)\,\underbrace{\sum_m w_m\,m^p}_{\,M_p\,}.
\]

By symmetry \(M_p = 0\) for \textbf{odd} \(p\) --- the odd Taylor terms
cancel automatically. The remaining moment conditions for an order-4
approximation of \(\partial^2 f / \partial x^2\) are:

{\def\LTcaptype{none} % do not increment counter
\begin{longtable}[]{@{}
  >{\raggedright\arraybackslash}p{(\linewidth - 4\tabcolsep) * \real{0.1111}}
  >{\raggedright\arraybackslash}p{(\linewidth - 4\tabcolsep) * \real{0.2222}}
  >{\raggedright\arraybackslash}p{(\linewidth - 4\tabcolsep) * \real{0.6667}}@{}}
\toprule\noalign{}
\begin{minipage}[b]{\linewidth}\raggedright
moment
\end{minipage} & \begin{minipage}[b]{\linewidth}\raggedright
required value
\end{minipage} & \begin{minipage}[b]{\linewidth}\raggedright
reason
\end{minipage} \\
\midrule\noalign{}
\endhead
\bottomrule\noalign{}
\endlastfoot
\(M_0 = w_0 + 2 w_1 + 2 w_2\) & \passthrough{\lstinline!0!} & kill the
constant (no \(f\) leak) \\
\(M_2 = 2 w_1 + 8 w_2\) & \passthrough{\lstinline!2!} & match
\(h^2 \cdot f''(x_i) \cdot (1/2)\) → coefficient
\passthrough{\lstinline!2!} after multiplying by
\passthrough{\lstinline"2!/h²"} \\
\(M_4 = 2 w_1 + 32 w_2\) & \passthrough{\lstinline!0!} & kill the
\(h^4\) term to lift accuracy from \(O(h^2)\) to \(O(h^4)\) \\
\end{longtable}
}

Solving:

\begin{itemize}
\tightlist
\item
  From \(M_2\) and \(M_4\): subtract → \(24 w_2 = -2\), so
  \(w_2 = -1/12\).
\item
  Back into \(M_2\): \(2 w_1 + 8 \cdot (-1/12) = 2\) →
  \(w_1 = (2 + 2/3)/2 = 4/3\).
\item
  From \(M_0\):
  \(w_0 = -2(w_1 + w_2) = -2(4/3 - 1/12) = -2 \cdot 5/4 \cdot ... = -5/2\).
  (Direct:
  \(w_0 = -2 \cdot 4/3 - 2 \cdot (-1/12) = -8/3 + 1/6 = -16/6 + 1/6 = -15/6 = -5/2\).)
\end{itemize}

Hence

\[
\boxed{\;w = \bigl(-\tfrac{1}{12},\ \tfrac{4}{3},\ -\tfrac{5}{2},\ \tfrac{4}{3},\ -\tfrac{1}{12}\bigr)\;}
\]

with truncation error from the next non-zero even moment
\(M_6 = 2 w_1 + 128 w_2 = 8/3 - 32/3 = -8\), giving leading error
\(-\tfrac{h^4}{6!}\cdot M_6 \cdot f^{(6)}(x_i) = +\tfrac{8 h^4}{720}\,f^{(6)}(x_i) = \tfrac{h^4}{90}\,f^{(6)}(x_i)\).

So the local truncation error is
\(\tfrac{h^4}{90}\,f^{(6)}(x_i) = O(h^4)\), which is the \textbf{order-4
accuracy} advertised. Compare with the order-2 stencil's
\(\tfrac{h^2}{12}\,f^{(4)}\).

\subsubsection{Order-4 first-derivative
stencil}\label{order-4-first-derivative-stencil}

Same procedure for \(\partial f / \partial x\) at \(x_i\):

\[
\partial_x f(x_i) \;\approx\; \frac{1}{h}\,\sum_{m=-2}^{2} w_m\,f(x_{i+m}),
\]

with \textbf{anti-symmetry} \(w_{-m} = -w_m\) (so \(w_0 = 0\),
\(w_{-1} = -w_1\), \(w_{-2} = -w_2\)). The non-trivial moment conditions
are now odd \(p\):

{\def\LTcaptype{none} % do not increment counter
\begin{longtable}[]{@{}
  >{\raggedright\arraybackslash}p{(\linewidth - 4\tabcolsep) * \real{0.1143}}
  >{\raggedright\arraybackslash}p{(\linewidth - 4\tabcolsep) * \real{0.2286}}
  >{\raggedright\arraybackslash}p{(\linewidth - 4\tabcolsep) * \real{0.6571}}@{}}
\toprule\noalign{}
\begin{minipage}[b]{\linewidth}\raggedright
moment
\end{minipage} & \begin{minipage}[b]{\linewidth}\raggedright
required value
\end{minipage} & \begin{minipage}[b]{\linewidth}\raggedright
reason
\end{minipage} \\
\midrule\noalign{}
\endhead
\bottomrule\noalign{}
\endlastfoot
\(M_1 = 2 w_1 + 4 w_2\) & \passthrough{\lstinline!1!} & match
\(h \cdot f'(x_i)\) → after \(1/h\) scale \\
\(M_3 = 2 w_1 + 16 w_2\) & \passthrough{\lstinline!0!} & kill the
\(h^3\) term to lift accuracy \\
\end{longtable}
}

Solving:

\begin{itemize}
\tightlist
\item
  Subtracting \(M_3 - M_1\): \(12 w_2 = -1\) → \(w_2 = -1/12\).
\item
  Back into \(M_1\): \(2 w_1 + 4 \cdot (-1/12) = 1\) →
  \(w_1 = (1 + 1/3)/2 = 2/3\).
\end{itemize}

Hence

\[
\boxed{\;w = \bigl(\tfrac{1}{12},\ -\tfrac{2}{3},\ 0,\ \tfrac{2}{3},\ -\tfrac{1}{12}\bigr)\;}
\]

Truncation error from \(M_5\):
\(M_5 = 2 w_1 + 64 w_2 = 4/3 - 16/3 = -4\), giving leading error
\(-\tfrac{h^4}{5!}\cdot M_5 \cdot f^{(5)}(x_i) = \tfrac{4 h^4}{120}\,f^{(5)}(x_i) = \tfrac{h^4}{30}\,f^{(5)}(x_i) = O(h^4)\).

Compare with the order-2 first-derivative stencil's
\(\tfrac{h^2}{6}\,f^{(3)}\).

\subsubsection{Why the per-axis-independent structure works on
anisotropic
grids}\label{why-the-per-axis-independent-structure-works-on-anisotropic-grids}

The two derivations above produce \textbf{unitless weight tables} ---
\(\{-1/12, 4/3, -5/2, 4/3, -1/12\}\) for \(\partial^2/\partial x^2\) and
\(\{1/12, -2/3, 0, 2/3, -1/12\}\) for \(\partial/\partial x\). Each
table encodes the linear combination of \emph{index-shifted samples}
that approximates the derivative; the spacing \(h\) enters only as a
scalar postmultiplication (\(1/h^2\) or \(1/h\)).

For 3D on a Cartesian-orthogonal grid, the Laplacian decomposes
axis-by-axis because the operator
\(\nabla^2 = \partial_x^2 + \partial_y^2 + \partial_z^2\) has no
cross-derivative terms --- the basis vectors are mutually perpendicular.
So we apply the same 1D weight table along each axis independently,
scaled by that axis's own cell size. The tables don't need to be
regenerated per axis.

The order-2 Phase 1 Laplacian already used this structure; Phase 7 just
substitutes a longer table.

\subsection{3. Algorithm}\label{algorithm}

The implementation is the smallest change consistent with the design:

\subsubsection{Stencil tables}\label{stencil-tables}

\href{../../../include/gridcalc/stencil/CentralDifference.hpp}{\passthrough{\lstinline!stencil/CentralDifference.hpp!}}
adds the explicit \passthrough{\lstinline!Coefficients<4>!}
specialization with \passthrough{\lstinline!radius = 2!},
\passthrough{\lstinline!offsets = \{-2, -1, 0, 1, 2\}!}, and the weights
derived above.
\href{../../../include/gridcalc/stencil/FirstDerivative.hpp}{\passthrough{\lstinline!stencil/FirstDerivative.hpp!}}
adds the analogous \passthrough{\lstinline!FirstDerivative<4>!}
specialization. Both follow the existing primary-undefined pattern:
\passthrough{\lstinline!Coefficients<3>!} /
\passthrough{\lstinline!FirstDerivative<3>!} etc. are still undefined
and trip a compile error at the call site.

\subsubsection{Operator templating}\label{operator-templating}

The three Phase 1/2 operator headers gain a
\passthrough{\lstinline!template <int Order = 2>!} parameter on the
function. The body changes by exactly one line --- the hardcoded stencil
type becomes the parameterized one:

\begin{lstlisting}
- using Coeffs = stencil::Coefficients<2>;
+ using Coeffs = stencil::Coefficients<Order>;
\end{lstlisting}

(and analogously for \passthrough{\lstinline!FirstDerivative<Order>!} in
\passthrough{\lstinline!Gradient.hpp!} and
\passthrough{\lstinline!Divergence.hpp!}). The per-axis loop already
iterates \passthrough{\lstinline!m!} over
\passthrough{\lstinline!Coeffs::offsets.size()!}, so a 5-element table
just runs the loop one more iteration per axis without any structural
change.

\subsubsection{Backward compatibility}\label{backward-compatibility}

C++17 resolves a function-template call with all template parameters
defaulted using the no-template-arg syntax:
\passthrough{\lstinline!laplacian(field)!} calls
\passthrough{\lstinline!laplacian<2>(field)!} because
\passthrough{\lstinline!Order!} defaults to \passthrough{\lstinline!2!}.
So every Phase 1/2 caller (and every Phase 4--6 transitive caller, like
\passthrough{\lstinline!func::evaluate(... f(\_, \_, lap\_psi))!} and
\passthrough{\lstinline!solve::diffuse!}) continues to compile and
produce bit-identical output. The Phase 1/2 test files are untouched;
the Phase 7 test file
(\passthrough{\lstinline!test/higher\_order\_test.cpp!}) exercises the
new \passthrough{\lstinline!<4>!} overload.

\subsubsection{Per-point cost}\label{per-point-cost}

Order 4 reads \passthrough{\lstinline!2 * radius + 1 = 5!} samples per
axis instead of \passthrough{\lstinline!3!}. Per-grid-point cost is
therefore \passthrough{\lstinline!\~5/3 ≈ 1.67×!} the order-2 cost ---
in exchange for the truncation error dropping from \(O(h^2)\) to
\(O(h^4)\), which at the typical test resolution
\passthrough{\lstinline!N = 32!} corresponds to \textbf{10--100×
absolute-accuracy improvement} (verified by
\passthrough{\lstinline!HigherOrderTest.*MoreAccurateThanOrder2!}). For
most applications this is a win; the Phase 20 performance pass will
quantify exact slowdown on the benchmark suite.

\subsection{4. Design decisions}\label{design-decisions}

\begin{itemize}
\tightlist
\item
  \textbf{Compile-time \passthrough{\lstinline!<Order>!} over runtime
  parameter.} Matches the project's existing pattern
  (\passthrough{\lstinline!stencil::Coefficients<Order>!} is already
  templated this way; Phase 3's
  \passthrough{\lstinline!Pairwise!}/\passthrough{\lstinline!Kahan!},
  Phase 4's arity dispatch, Phase 6's
  \passthrough{\lstinline!ExplicitEuler!}/\passthrough{\lstinline!RK4!}).
  No runtime branch in the inner loop. Compile-time error on unsupported
  orders.
\item
  \textbf{Hard-coded weights, derivation in the Developer Note.} Two
  specializations, \textasciitilde10 LOC each. Deferred constexpr
  Fornberg until Phase 8/11 actually require many more orders. See
  \passthrough{\lstinline!requirements.md!} for the full rationale.
\item
  \textbf{Default \passthrough{\lstinline!Order = 2!} rather than
  removing the default.} Backward compatibility --- Phase 1/2/4/5/6
  callers don't need to change. C++17 function-template default-argument
  semantics make this clean.
\item
  \textbf{Apply to all three operators in one phase.} Matches the
  roadmap deliverable list. Divergence is templated for symmetry even
  though the roadmap text only names Laplacian and gradient --- the
  operators are tightly coupled (divergence reads gradient's
  \passthrough{\lstinline!FirstDerivative!} table) and skipping
  divergence would force a Phase 7.5 follow-up.
\item
  \textbf{No \passthrough{\lstinline!Order!} parameter on
  \passthrough{\lstinline!solve::diffuse!}.}
  \passthrough{\lstinline!solve::diffuse!} is the thin convenience
  driver for \(D \nabla^2 \psi\); it currently calls
  \passthrough{\lstinline!diff::laplacian(...)!} without a template arg,
  which now resolves to order 2. A future phase can promote the
  diffusion driver to take an \passthrough{\lstinline!Order!} parameter
  if the use case justifies it; today the workaround (call
  \passthrough{\lstinline!solve::integrate!} directly with a custom RHS
  that uses \passthrough{\lstinline!laplacian<4>!}) is documented in the
  User Guide.
\end{itemize}

\subsection{5. References}\label{references}

{[}1{]} Fornberg, B. (1988). ``Generation of Finite Difference Formulas
on Arbitrarily Spaced Grids.'' \emph{Mathematics of Computation},
51(184), 699--706.
\url{https://doi.org/10.1090/S0025-5718-1988-0935077-0}. The canonical
reference for deriving central-difference weights at arbitrary order on
arbitrary stencils. Used here as the \emph{theoretical} derivation
framework even though the two specific specializations are hard-coded;
the derivation in §2 above is a hand-evaluation of Fornberg's algorithm
at order 4 with stencil offsets
\passthrough{\lstinline!\{-2, -1, 0, 1, 2\}!}.

{[}2{]} LeVeque, R. J. (2007). \emph{Finite Difference Methods for
Ordinary and Partial Differential Equations}. SIAM. ISBN
978-0-89871-783-9. \url{https://doi.org/10.1137/1.9780898717839}. §1.5
covers higher-order finite differences and lists the same
\passthrough{\lstinline!\{-1/12, 4/3, -5/2, 4/3, -1/12\}!} stencil for
\(\partial^2/\partial x^2\) in Table 1.1.

{[}3{]} DLMF (Digital Library of Mathematical Functions), §3.4
``Differentiation.'' \url{https://dlmf.nist.gov/3.4} (permanent URL).
NIST's online reference includes the central-difference weights for both
first and second derivatives at orders 2, 4, 6, 8 in tabular form;
matches the derivation here.

{[}4{]} Trefethen, L. N. (2000). \emph{Spectral Methods in MATLAB}.
SIAM. \url{https://doi.org/10.1137/1.9780898719598}. Chapter 1 motivates
the choice between increasing stencil order vs refining the grid (the
classic FD-vs-spectral tradeoff); the Phase 7 User Guide note's ``When
to pick order 4'' guidance is in the spirit of this discussion.
